\documentclass[]{article}
% \usepackage[utf8]{inputenc}
\usepackage[polish]{babel}
\usepackage{amsmath}
\usepackage{amssymb}
\usepackage{mathtools}
\usepackage{geometry}
\usepackage{fancyhdr}
\usepackage{graphicx}
\usepackage{xcolor}
\usepackage{hyperref}

\geometry{margin=2.5cm}
\pagestyle{fancy}
\fancyhf{}
\fancyhead[R]{\thepage}
\fancyhead[L]{Raport Matematyczny ABM 2.0}

\title{\textbf{Raport Wzorów Matematycznych} \\
\textit{Model Symulacyjny Agent-Based Model (ABM 2.0)}}
\author{Projekt Zespołowy}
\date{\today}

\begin{document}

\maketitle

\tableofcontents
\newpage

\section{Wprowadzenie}

Niniejszy raport dokumentuje wszystkie wzory matematyczne wykorzystane w modelu symulacyjnym ABM 2.0. Model odwzorowuje dynamikę demograficzną populacji z uwzględnieniem urodzin, zgonów, powstawania/rozpadu gospodarstw domowych oraz wpływu chorób na życie obywateli.

\section{Ryzyka Śmiertelności}

\subsection{Ogólny Model Ryzyka Śmiertelności}

Ryzyko śmiertelności dla każdego obywatela obliczane jest miesięcznie na podstawie czterech głównych komponentów: wieku, chorób, niepełnosprawności oraz płci.

\subsubsection{Czynnik Wieku}

Czynnik uwarunkowany wiekiem obliczany jest za pomocą funkcji wykładniczej:

\begin{equation}
R_{age} = \lambda_0 \cdot \left(1 + e^{\frac{a - a_0}{\sigma_a}}\right)
\end{equation}

gdzie:
\begin{itemize}
    \item $\lambda_0 = 0.00005$ --- bazowa miesięczna stopa śmiertelności
    \item $a$ --- wiek osoby (w latach)
    \item $a_0 = 55$ --- wiek referencyjny (punkt przegięcia)
    \item $\sigma_a = 15$ --- parametr skalowania
\end{itemize}

\subsubsection{Czynnik Chorób i Niepełnosprawności}

Czynnik związany z chorobami i niepełnosprawością jest obliczany jako kombinacja liniowa liczby chorób oraz wskaźnika niepełnosprawności:

\begin{equation}
R_{disease} = \mu \left(0.001 \cdot n_{cond} + 0.005 \cdot d_i\right)
\end{equation}

gdzie:
\begin{itemize}
    \item $\mu = 1.0$ --- mnożnik chorób (parametr kalibracji modelu)
    \item $n_{cond}$ --- liczba chorób (warunków zdrowotnych)
    \item $d_i$ --- wskaźnik niepełnosprawności osoby
    \item współczynniki 0.001 i 0.005 --- wagi dla liczby chorób i niepełnosprawności
\end{itemize}

\subsubsection{Całkowite Ryzyko Śmiertelności}

Całkowite miesięczne ryzyko śmiertelności obliczane jest jako suma komponentów z uwzględnieniem przewagi kobiet:

\begin{equation}
R_{mortality}(t) = \begin{cases}
(R_{age} + R_{disease}) \cdot f_{sex} & \text{jeśli } sex = \text{kobieta} \\
R_{age} + R_{disease} & \text{jeśli } sex = \text{mężczyzna}
\end{cases}
\end{equation}

gdzie $f_{sex} = 0.85$ --- współczynnik redukcji śmiertelności dla kobiet (kobiety mają 15\% niższe ryzyko).

Ostateczne ryzyko jest ograniczane do przedziału $[0, 1]$:

\begin{equation}
R_{final} = \min\left(\max(R_{mortality}, 0), 1\right)
\end{equation}

\subsubsection{Decyzja o Śmierci}

Śmierć danej osoby w miesiącu $t$ jest modelowana jako zdarzenie bernoullowskie:

\begin{equation}
Death_i(t) = \begin{cases}
True & \text{jeśli } U < R_{final} \\
False & \text{w pozostałych przypadkach}
\end{cases}
\end{equation}

gdzie $U \sim \text{Uniform}(0,1)$ --- zmienna losowa o rozkładzie jednostajnym.

\section{Model Płodności i Narodzin}

\subsection{Prawdopodobieństwo Urodzenia}

Prawdopodobieństwo urodzenia dziecka w danym miesiącu obliczane jest tylko dla kobiet w wieku reprodukcyjnym (18-40 lat). Model uwzględnia wiek, okres szczytowej płodności oraz wpływ chorób.

\subsubsection{Bazowa Funkcja Płodności}

Bazowe ryzyko płodności modelowane jest za pomocą rozkładu normalnego (Gaussian) skoncentrowanego na wieku szczytowej płodności (27,5 lat):

\begin{equation}
f_{base}(a) = e^{-\frac{(a - a_{peak})^2}{\sigma^2}}
\end{equation}

gdzie:
\begin{itemize}
    \item $a$ --- wiek kobiety (w latach)
    \item $a_{peak} = 27.5$ --- wiek szczytowej płodności
    \item $\sigma^2 = 25$ --- parametr wariancji rozkładu normalnego
\end{itemize}

\subsubsection{Miesięczna Stopa Płodności}

Miesięczna stopa płodności obliczana jest jako:

\begin{equation}
r_{fertility}(a) = \begin{cases}
f_{base}(a) \cdot 0.003 & \text{jeśli } 18 \le a \le 40 \\
0 & \text{w pozostałych przypadkach}
\end{cases}
\end{equation}

gdzie współczynnik 0.003 odpowiada w przybliżeniu 3.6\% rocznej stopie płodności na szczycie.

\section{Płodność Kobiet: Wiek i Determinanty Medyczne}

Płodność kobiet jest kluczowym aspektem modelu demograficznego. Dane medyczne wskazują, że szczyt płodności przypada na wiek 20–25 lat, kiedy organizm osiąga najwyższą sprawność hormonalną i genetyczną. Po 30. roku życia płodność zaczyna stopniowo spadać, a po 35. roku życia proces ten przyspiesza.

\subsection{Okno Optymalnego Wieku Prokreacyjnego}

Statystyczne prawdopodobieństwo zajścia w ciążę w jednym cyklu miesięcznym w zależności od wieku przedstawia się następująco:

\begin{table}[h!]
\centering
\begin{tabular}{|c|c|c|}
\hline
\textbf{Grupa wiekowa (lata)} & \textbf{Szansa na ciążę w cyklu (\\%)} & \textbf{Prawdopodobieństwo ciąży po 12 cyklach (")} \\
\hline
19–25 & 25\% & ~92\% \\
26–29 & 20\% & ~85–90\% \\
30–34 & 15–20\% & ~75–80\% \\
35–39 & 10–15\% & ~50–60\% \\
40–44 & 5–7\% & ~25–30\% \\
45+ & <1–5\% & <10\% \\
\hline
\end{tabular}
\caption{Prawdopodobieństwo zajścia w ciążę w zależności od wieku.}
\end{table}

\subsection{Ryzyko Poronień}

Ryzyko utraty ciąży wzrasta wraz z wiekiem, głównie z powodu aberracji chromosomalnych. Tabela poniżej przedstawia ryzyko poronienia w różnych grupach wiekowych:

\begin{table}[h!]
\centering
\begin{tabular}{|c|c|}
\hline
\textbf{Wiek kobiety (lata)} & \textbf{Ryzyko poronienia (")} \\
\hline
<20 & 17\% \\
20–24 & 10–11\% \\
25–29 & 10\% \\
30–34 & 11–12\% \\
35–39 & 17–25\% \\
40–44 & 33–40\% \\
45+ & 57–80\% \\
\hline
\end{tabular}
\caption{Ryzyko poronienia w zależności od wieku.}
\end{table}

\subsection{Fizjologiczne Podłoże Płodności}

Płodność kobiet zależy od rezerwy jajnikowej, która zmniejsza się wraz z wiekiem. W wieku 20–25 lat rezerwa jajnikowa jest najwyższa, co zapewnia optymalną jakość oocytów. Po 35. roku życia liczba pęcherzyków jajnikowych gwałtownie spada, co prowadzi do zmniejszenia potencjału rozrodczego.

\section{Uwagi i Ograniczenia}

\begin{itemize}
    \item Model wykorzystuje podejście stochastyczne z generatorem liczb losowych
    \item Wszystkie prawdopodobieństwa są ograniczane do przedziału $[0, 1]$
    \item Parametry mogą być kalibrowane na podstawie danych rzeczywistych
    \item Model zakłada niezależność zdarzeń demograficznych w ramach jednego miesiąca
    \item Choroby są modelowane binarnie (osoba ma lub nie ma chorobę)
\end{itemize}

\end{document}